%====================================================================
%  Tidal Lanes — Estimation Notes
%
%  Companion to TidalLanes_Model_v4.tex.  Structure follows
%  Path_Dependence_in_Transportation_Network_Expansion_estimation_notes.pdf.
%  Last updated: 2026-06-30 (gate5 mainline bpr+pos05).
%====================================================================
\documentclass[11pt]{article}

\usepackage[margin=1in]{geometry}
\usepackage{amsmath,amssymb,amsthm,mathtools}
\usepackage{booktabs,longtable,enumitem}
\usepackage[colorlinks=true,linkcolor=blue,citecolor=blue,urlcolor=blue]{hyperref}
\usepackage{setspace}
\newtheorem{assumption}{Assumption}
\setstretch{1.12}
\setlength{\parindent}{0pt}
\setlength{\parskip}{0.45em}

\newcommand{\Ncal}{\mathcal{N}}
\newcommand{\Lbar}{\bar L}
\newcommand{\Wbar}{\bar W}
\newcommand{\ubar}{\bar u}
\newcommand{\Abar}{\bar A}
\newcommand{\obs}{\mathrm{obs}}
\newcommand{\phys}{\mathrm{phys}}
\newcommand{\AAcal}{\mathrm{AA}}

\title{Tidal Lanes --- Estimation Notes}
\author{Working draft (gate5 mainline)}
\date{June 30, 2026}

\begin{document}
\maketitle

%====================================================================
\section{Overview}\label{sec:overview}
%====================================================================

This note records how each object in the directional spatial commuting
model is \emph{measured}, \emph{calibrated}, \emph{inverted}, or
\emph{estimated}.  Equation labels refer to
\texttt{TidalLanes\_Model\_v4.tex} unless noted otherwise.

\paragraph{Parameter accounting.}
We distinguish three tiers:
\begin{enumerate}[leftmargin=2em]
  \item \textbf{Data construction} (no free parameters): directed
        network primitives, peak-period travel times, lane counts, and
        road-mode OD flows.
  \item \textbf{Offline calibration} (fixed from external literature or
        documented mapping exercises): $(\theta,\alpha,\beta)$ and the
        congestion elasticities $(\lambda_{\phys},\lambda_{\AAcal})$.
  \item \textbf{Inversion} (recovers latent fundamentals given
        Assumption~\ref{ass:obs_eq}): $(\ubar_i,\Abar_j,\chi)$.
  \item \textbf{Joint / structural estimation} (not yet implemented for
        the gate5 mainline): panel or IV congestion elasticity; optional
        GMM on sorting parameters.
\end{enumerate}

\paragraph{Status (2026-06-30).}
Sections~\ref{sec:data}--\ref{sec:invert} are \emph{complete} for the
gate5 mainline bundle
(\texttt{model\_input\_edgeproj\_v2\_lcc\_gate5\_v1}).
Section~\ref{sec:diagnostic} reports \emph{diagnostic} regressions only.
Section~\ref{sec:planned} lists open structural estimation tasks.

\paragraph{Mainline counterfactual specification (for reference).}
Pair selection \texttt{score\_bpr\_st}; $\alpha=+0.05$, $\beta=-0.10$;
anchored congestion with $\lambda_{\phys}=1.1$; tidal policy top-100
pairs, $\Delta n=1$.  Full numerical results:
\texttt{Documents/L2\_notes/model\_and\_results\_summary\_20260630.md}.

%====================================================================
\section{Data Construction}\label{sec:data}
%====================================================================

These objects enter the model as \emph{primitives} (v4 \S6.1, Table~1).
They are not estimated within the structural system.

\begin{longtable}{@{}p{0.42\linewidth}p{0.18\linewidth}p{0.32\linewidth}@{}}
\toprule
\textbf{Description} & \textbf{Symbol / object} & \textbf{Source (gate5 mainline)} \\
\midrule
\endhead
Spatial cells (LCC) & $N=1104$ & Square grid, gate5 largest connected component \\
Directed edges & $|\mathcal{E}|\approx 2642$ & \texttt{qsm\_input\_edges\_square.parquet} \\
Directed length & $\ell_{kl}$ (km) & QSM stage-06 export \\
Free-flow speed / time & $v^0_{kl}$, $\bar\tau_{kl}$ & \texttt{tau\_ff\_min} \\
Observed peak speed / time & $v^{\obs}_{kl}$, $\tau^{\obs}_{kl}$ & Baidu AM peak; edgeproj\_v2 \\
Directed lane count & $n_{kl}$ & Lane lookup (2026-06-29 gate5 lane-fix) \\
Bilateral OD flow & $L^{\obs}_{ij}$ & \texttt{commuters\_road} (default column) \\
Residence / workplace shares & $l^{R,\obs}_i$, $l^{F,\obs}_j$ & Marginals of $L^{\obs}_{ij}$ \\
Total commuters & $\Lbar$ & $\sum_i L^{R,\obs}_i \approx 88{,}110$ \\
\bottomrule
\end{longtable}

\paragraph{Implementation.}
\texttt{data\_work/src/model/full\_estimation/build\_primitives.py}
loads the parquet bundle under
\texttt{data\_work/outputs/<version\_id>/}, reindexes nodes to
contiguous $0..N-1$, and applies lane fallbacks documented in
\texttt{data\_work/docs/STRUCTURAL\_MODEL\_WORKFLOW.md}.

\paragraph{Open design choices (switchable, not locked).}
\begin{itemize}[leftmargin=2em]
  \item Grid geometry: square (baseline), hex, voronoi.
  \item Peak period: AM (mainline), PM robustness.
  \item Travel-cost pipeline: raw connected vs.\ edge projection v2
        (mainline uses \texttt{edgeproj\_v2\_lcc\_gate5}).
  \item OD column: \texttt{commuters\_road} vs.\ mode-specific columns.
\end{itemize}

%====================================================================
\section{Transportation Network and Directed Travel Costs}\label{sec:transport}
%====================================================================

\begin{longtable}{@{}p{0.42\linewidth}p{0.22\linewidth}p{0.28\linewidth}@{}}
\toprule
\textbf{Description} & \textbf{Current value} & \textbf{Status} \\
\midrule
\endhead
Median $\tau^{\obs}/\tau^{\mathrm{ff}}$ (edges) & $\approx 1.00--1.35$ & Measured \\
Share of edges with $\tau^{\obs}<\tau^{\mathrm{ff}}$ & $\approx 8.3\%$ & QA flag (anchored floor) \\
Baseline median $\tau/\tau^{\obs}$ (endogenous) & $1.0001$ & Converged (mainline) \\
Spectral-radius rescaling & $\kappa\ge 1$ (auto) & Algorithmic stabilizer \\
\bottomrule
\end{longtable}

\subsection{Iceberg commuting cost index}

From v4 \S3.2--3.3, the bilateral cost index satisfies
\begin{equation}\label{eq:tau_series}
  \tau_{ij}^{-\theta} = \bigl[(I-A)^{-1}\bigr]_{ij},
  \qquad
  A_{kl} = t_{kl}^{-\theta},\quad
  t_{kl} = (\tau^{\min}_{kl})^{-1/\theta}.
\end{equation}
At the data stage we set $\tau^{\min,\obs}_{kl}=\tau^{\obs}_{kl}$ and
build $\tau^{\obs}_{ij}(\theta)$ once per calibration of $\theta$.

\subsection{Congestion technology (three implemented forms)}

Physical minutes update (level form, v4 \S3.6):
\begin{equation}\label{eq:congestion_legacy}
  \tau^{\min}_{kl}
  \;=\;
  \bar\tau_{kl}\,
  \Bigl(\frac{\Xi_{kl}}{n_{kl}}\Bigr)^{\lambda_{\AAcal}}
  \qquad\text{(legacy / AA iceberg mapping).}
\end{equation}

Density form (Barwick units):
\begin{equation}\label{eq:congestion_density}
  \tau^{\min}_{kl}
  \;=\;
  \bar\tau_{kl}\,
  \Bigl(\frac{\Xi_{kl}}{n_{kl}\ell_{kl}}\Bigr)^{\lambda_{\phys}}.
\end{equation}

Anchored form (gate5 mainline):
\begin{equation}\label{eq:congestion_anchored}
  \tau^{\min}_{kl}
  \;=\;
  \max\Bigl\{
    \bar\tau_{kl},\;
    \tau^{\obs}_{kl}\,
    \Bigl(\frac{\Xi_{kl}}{\Xi^{\obs}_{kl}}\frac{n^{\obs}_{kl}}{n_{kl}}\Bigr)^{\lambda_{\phys}}
  \Bigr\}.
\end{equation}
Anchored congestion \emph{semi-locks} the baseline at observed travel
times when $\Xi=\Xi^{\obs}$ and $n=n^{\obs}$, while preserving the
congestion \emph{slope} $\lambda_{\phys}$ around that point.  This is
the default in \texttt{equilibrium\_solve.py}
(\texttt{MAINLINE\_CONGESTION\_FORM = "anchored"}).

\paragraph{Code.}
Congestion updates live in
\texttt{data\_work/src/model/full\_estimation/equilibrium\_solve.py};
observation anchors $(\Xi^{\obs}_{kl},n^{\obs}_{kl})$ in
\texttt{compute\_observation\_anchors()}.

%====================================================================
\section{Congestion Elasticity}\label{sec:lambda}
%====================================================================

\begin{longtable}{@{}p{0.38\linewidth}p{0.22\linewidth}p{0.34\linewidth}@{}}
\toprule
\textbf{Description} & \textbf{Symbol} & \textbf{Current value (mainline)} \\
\midrule
\endhead
Physical travel-time elasticity w.r.t.\ density &
  $\lambda_{\phys}$ & $\mathbf{1.1}$ (Barwick et al.\ 2024) \\
AA / iceberg-cost elasticity &
  $\lambda_{\AAcal}=\lambda_{\phys}/\theta$ & $\approx 0.161$ \\
Legacy AA replication import &
  $\delta_1$ & $0.488$ (Allen--Arkolakis 2022 Seattle) \\
Legacy code default &
  $\lambda_{\AAcal}=\delta_1/\theta$ & $\approx 0.0714$ \\
Robustness band (Phase A) &
  $\lambda_{\phys}$ & $0.572$ (legacy $\delta_1$ as physical slope) \\
\bottomrule
\end{longtable}

\subsection{External mapping (Barwick et al.\ 2024)}

Barwick, Li, and Lin estimate, for Beijing traffic sensors,
\begin{equation}
  \ln S_{st} = e\,\ln D_{st} + \cdots,
  \qquad \hat e \approx -1.1,
\end{equation}
where $S$ is speed and $D$ is traffic density (vehicles per lane-km).
Since $\tau_{kl}=\ell_{kl}/S_{kl}$ (holding length fixed),
\begin{equation}\label{eq:lambda_phys_barwick}
  \lambda_{\phys}
  \;=\;
  -\frac{\partial\ln S_{kl}}{\partial\ln D_{kl}}
  \;\approx\;
  1.1.
\end{equation}
Under the AA iceberg mapping $t_{kl}=\tau_{kl}^{1/\theta}$,
\begin{equation}\label{eq:lambda_AA_barwick}
  \lambda_{\AAcal} = \frac{\lambda_{\phys}}{\theta} \approx \frac{1.1}{6.83} \approx 0.161.
\end{equation}
Full derivation:
\texttt{data\_work/docs/CONGESTION\_ELASTICITY\_PARAMETER\_MAPPING.tex}.

\subsection{AA Seattle replication import (legacy default)}

The replication scaffold imports Allen and Arkolakis (2022) Seattle
congestion coefficient $\delta_1=0.488$ and sets
\begin{equation}
  \lambda_{\AAcal} = \frac{\delta_1}{\theta} = \frac{0.488}{6.83} \approx 0.0714,
\end{equation}
stored in \texttt{parameters.py} as \texttt{Parameters.lambda\_}.
This value governs the \texttt{legacy} congestion form
\eqref{eq:congestion_legacy}.  It is \textbf{not} the mainline anchored
slope; mainline counterfactuals use $\lambda_{\phys}=1.1$ in
\eqref{eq:congestion_anchored}.

\subsection{Why structural $\lambda$ is not yet identified}

v4 \S6.3 states two data
limitations:
\begin{enumerate}[leftmargin=2em]
  \item Directed link traffic $\Xi_{kl}$ is \emph{not observed} in the
        QSM bundle; only speeds, free-flow speeds, lanes, and OD flows.
  \item We observe a single peak-period cross-section, not a panel with
        exogenous flow variation suitable for IV.
\end{enumerate}
Until Baidu crawl / sensor panel data are merged at the directed-edge
level, $\lambda_{\phys}$ is \emph{calibrated} from external evidence
(Barwick) rather than structurally estimated.

%====================================================================
\section{Sorting and Commuting Parameters}\label{sec:sorting}
%====================================================================

\begin{longtable}{@{}p{0.38\linewidth}p{0.18\linewidth}p{0.36\linewidth}@{}}
\toprule
\textbf{Description} & \textbf{Symbol} & \textbf{Current value} \\
\midrule
\endhead
Fr\'echet shape (commuting / route) & $\theta$ & $\mathbf{6.83}$ \\
Production agglomeration elasticity & $\alpha$ & $\mathbf{+0.05}$ (mainline) \\
Residential amenity elasticity & $\beta$ & $\mathbf{-0.10}$ \\
\bottomrule
\end{longtable}

\subsection{Fr\'echet shape $\theta$}

\paragraph{Baseline calibration.}
$\theta=6.83$ is taken from Allen and Arkolakis (2022) Seattle
replication, which follows Ahlfeldt et al.\ (2015).  It governs:
\begin{itemize}[leftmargin=2em]
  \item route choice within the Fr\'echet framework (v4 \S3.2);
  \item the commuting gravity exponent in
        $\ln L_{ij} \propto -\theta\ln\tau_{ij}$;
  \item the mapping $\lambda_{\AAcal}=\lambda_{\phys}/\theta$.
\end{itemize}

\paragraph{Internal diagnostic (not used to override baseline).}
v4 \S6.5 proposes a grid search over
$\theta\in\{3,\ldots,12\}$: for each candidate, recompute
$\tau^{\obs}_{ij}(\theta)$ from \eqref{eq:tau_series} and estimate
\begin{equation}\label{eq:gravity_log}
  \ln L^{\obs}_{ij}
  \;=\;
  -\theta\,\ln\tau^{\obs}_{ij}(\theta) + \mu_i + \nu_j + \varepsilon_{ij},
\end{equation}
selecting the $\theta$ at which the regression slope equals the imposed
$\theta$.  On early square-grid prototypes this diagnostic yielded
$\hat\theta$ far below $6.83$ (e.g.\ $0.62$--$0.70$); the external
calibration remains baseline because $\tau^{\obs}_{ij}(\theta)$ is itself
an equilibrium object.  Outputs archived under
\texttt{data\_work/outputs/match\_projection\_complete\_v7\_latest/}.

\subsection{Agglomeration elasticities $(\alpha,\beta)$}

\paragraph{Sign convention.}
The model notes write $\alpha>0$ and $\beta<0$ for the \emph{economic}
signs of production spillovers and residential congestion.  The AA
replication code uses iceberg-cost elasticities with \emph{negative}
numeric values ($\alpha=-0.12$, $\beta=-0.10$).  The level-form solver
uses the replication signs in \texttt{parameters.py} defaults; the
gate5 \textbf{mainline} overrides to gross-spillover calibration
$\alpha=+0.05$, $\beta=-0.10$ (see \texttt{run\_stage2\_gate5.py} and
\texttt{stage2\_.../parameters.json}).

\paragraph{Sources.}
\begin{itemize}[leftmargin=2em]
  \item $\beta=-0.10$: held from AA (2022) Seattle.
  \item $\alpha=+0.05$: mid-range gross production spillover from the
        urban / spatial literature (mainline ``pos05'' track); replaces
        the replication default $-0.12$ for counterfactual analysis.
\end{itemize}

\paragraph{Sensitivity (converged grid, gate5 bpr).}
With $\beta=-0.10$ fixed, person-min/day crosses zero near
$\alpha\in[0.04,0.05]$, while $\Wbar_{\mathrm{LT}}$ remains negative
throughout the scanned region (best: $\alpha=-0.12$, $\Wbar_{\mathrm{LT}}\approx -0.32\%$).
Full grid:
\texttt{outputs/alpha\_beta\_sensitivity\_bpr\_pos05\_v1/}.
\textbf{No converged specification yields $\Wbar_{\mathrm{LT}}>0$.}

\paragraph{Estimation status.}
$(\theta,\alpha,\beta)$ are \textbf{not} jointly estimated in the
current pipeline.  Planned extensions: sensitivity tables (done),
optional GMM moments on second moments of $(l^R,l^F)$ changes under
transport shocks (not implemented).

%====================================================================
\section{Inversion of Fundamentals}\label{sec:invert}
%====================================================================

\begin{longtable}{@{}p{0.42\linewidth}p{0.48\linewidth}@{}}
\toprule
\textbf{Target object} & \textbf{Status} \\
\midrule
\endhead
$(\ubar_i^{\theta},\Abar_j^{\theta})$ & Recovered (gate5 mainline) \\
Scalar $\chi$ & Recovered jointly \\
Normalization & $\mathrm{gmean}(\ubar_i^{\theta})=\mathrm{gmean}(\Abar_j^{\theta})=1$ \\
\bottomrule
\end{longtable}

\subsection{Assumption and equations}

\begin{assumption}[Observed allocation is an equilibrium]\label{ass:obs_eq}
$(L_i^{R,\obs},L_j^{F,\obs})$ equals the equilibrium allocation in
v4 Definition~4.1 under observed travel costs, calibrated
$(\theta,\alpha,\beta,\lambda)$, and some baseline fundamentals
$(\ubar_i,\Abar_j)$.
\end{assumption}

Under Assumption~\ref{ass:obs_eq}, v4 eqs.\ (invert\_R)--(invert\_F):
\begin{align}
  (l_i^{R,\obs})^{1-\beta\theta}
  &= \chi\sum_j (\tau^{\obs}_{ij})^{-\theta}\,
     \ubar_i^{\theta}\,\Abar_j^{\theta}\,(l_j^{F,\obs})^{\alpha\theta},
  \label{eq:invert_R}\\
  (l_j^{F,\obs})^{1-\alpha\theta}
  &= \chi\sum_i (\tau^{\obs}_{ij})^{-\theta}\,
     \ubar_i^{\theta}\,\Abar_j^{\theta}\,(l_i^{R,\obs})^{\beta\theta}.
  \label{eq:invert_F}
\end{align}

\subsection{Algorithm}

\texttt{invert\_fundamentals.py} alternates updates on
$v_i=\ubar_i^{\theta}$ and $w_j=\Abar_j^{\theta}$ with geometric-mean
normalization each step; recomputes $\chi$ from the implied equation
constants; stops when $\max|\Delta\log v|,|\Delta\log w|<\texttt{tol}$.

\paragraph{Implementation details.}
\begin{itemize}[leftmargin=2em]
  \item $(I-A^{\obs})^{-1}$ computed once with dense linear algebra
        ($N=1104$).
  \item Floor $10^{-6}$ on zero shares before normalization.
  \item Residuals \texttt{resid\_R}, \texttt{resid\_F} logged; large
        values flag mismatch between data shares and assumed equilibrium.
\end{itemize}

\paragraph{Outputs (gate5 mainline).}
\texttt{stage2\_.../fundamentals/} $\rightarrow$
\texttt{ubar\_theta}, \texttt{abar\_theta}, \texttt{chi}, inversion
meta JSON.  Fundamentals are \textbf{held fixed} in all counterfactuals.

\paragraph{Diagnostic maps.}
\texttt{run\_fundamental\_maps\_gate5.py} reports
$\mathrm{corr}(\ubar_i^{\theta},L_i^{R,\obs})\approx 0.30$ and
$\mathrm{corr}(\Abar_j^{\theta},L_j^{F,\obs})\approx 0.67$ on gate5.

%====================================================================
\section{Diagnostic and Reduced-Form Checks}\label{sec:diagnostic}
%====================================================================

These routines \textbf{do not} overwrite baseline parameters.  They
provide calibration checks and functional-form guidance.

\subsection{Cross-section congestion regression (model-implied $\hat\Xi$)}

Target (v4 congestion regression):
\begin{equation}\label{eq:lambda_regression}
  \ln\!\left(\frac{\tau^{\obs}_{kl}}{\bar\tau_{kl}}\right)
  \;=\;
  \lambda_{\AAcal}\,
  \ln\!\left(\frac{\hat\Xi_{kl}}{n_{kl}}\right)
  + \eta_{kl} + \varepsilon_{kl}.
\end{equation}

\paragraph{Procedure.}
\texttt{lambda\_diagnostic.estimate\_lambda\_cross\_section}:
\begin{enumerate}[leftmargin=2em]
  \item Obtain $\hat\Xi_{kl}$ from baseline equilibrium at observed
        inputs.
  \item Demean $\ln(\tau^{\obs}/\bar\tau)$ and $\ln(\hat\Xi/n)$ by
        integer lane bins (proxy for edge fixed effects).
  \item OLS slope $\Rightarrow$ \texttt{lambda\_hat}.
\end{enumerate}

\paragraph{Gate5 mainline result}
(\texttt{stage2\_.../lambda\_diagnostic.json}):
\begin{center}
\begin{tabular}{@{}lr@{}}
\toprule
$\hat\lambda_{\AAcal}$ & $0.00469$ \\
Heteroskedasticity-robust s.e. & $(0.00183)$ \\
$R^2$ & $0.0025$ \\
$N$ & $2642$ \\
\bottomrule
\end{tabular}
\end{center}

\paragraph{Interpretation.}
The estimate is an order of magnitude below the AA import ($0.071$) and
two orders below the Barwick-implied $\lambda_{\AAcal}\approx 0.161$.
This is expected: $\hat\Xi_{kl}$ is circular (model-assigned, not
measured), and a single cross-section lacks credible instruments.
\textbf{Not used} to set mainline $\lambda_{\phys}=1.1$.

\subsection{Reduced-form congestion horse race (P0)}

\texttt{run\_congestion\_reduced\_form\_gate5.py} compares parametric
families for $\ln(\tau^{\obs}/\tau^{\mathrm{ff}})$ as a function of
capacity index $z=1/(n_{kl}\ell_{kl})$, including BPR-style proxies
inferred \emph{only} from $(\tau^{\obs},\tau^{\mathrm{ff}})$ without
equilibrium traffic.  Outputs (when run):
\texttt{outputs/congestion\_reduced\_form\_gate5\_v1/}.

\subsection{Equilibrium convergence as implicit diagnostic}

Mainline endogenous baseline (gate5 bpr+pos05):
\begin{itemize}[leftmargin=2em]
  \item converged in $65$ outer iterations;
  \item $\max\log$ change $\approx 10^{-4}$;
  \item median $\tau/\tau^{\obs}=1.0001$.
\end{itemize}
Positive $\alpha=0.05$ requires damped solver
(\texttt{damping=0.10}, \texttt{max\_outer=500}) in sensitivity cells.

%====================================================================
\section{Equilibrium Solution and Counterfactuals}\label{sec:equilibrium}
%====================================================================

Given inverted fundamentals, the level-form solver
(\texttt{equilibrium\_solve.py}) nests:
\begin{enumerate}[leftmargin=2em]
  \item Outer loop on $\tau^{\min}_{kl}$ via \eqref{eq:congestion_anchored};
  \item Inner share fixed point $(l^R,l^F)$ (v4 eqs.\ eq\_R--eq\_F);
  \item Market-access fixed point $(\Pi,P)$ (v4 eq.\ E4);
  \item Flow assignment $\Xi_{kl}$ (v4 eq.\ E5).
\end{enumerate}

Counterfactuals re-solve with $(\ubar_i^{\theta},\Abar_j^{\theta})$
fixed.  Welfare:
\begin{equation}
  \Wbar^{\theta} = \sum_i u_i^{\theta}\,\Pi_i^{-\theta},
  \qquad
  u_i^{\theta} = \ubar_i^{\theta}(L_i^R)^{\beta\theta}.
\end{equation}

\paragraph{Reported decomposition (mainline top100).}
ST fixed-share gain $\approx +4.6$ pp; sorting total $\approx -8.5$ pp;
net $\Wbar_{\mathrm{LT}}\approx -0.85\%$; person-min/day $\approx +95$.
See welfare diagnostics CSV in \S\ref{sec:registry}.

%====================================================================
\section{Planned Structural Estimation}\label{sec:planned}
%====================================================================

Following the Path Dependence notes' split between offline calibration
and joint GMM, the TidalLanes pipeline \textbf{has not yet} implemented
a just-identified GMM stage.  Candidate moment blocks:

\subsection{Panel congestion elasticity (priority)}

When Baidu crawl / sensor panel delivers directed flows or densities
$\{D_{kl,t}\}$:
\begin{equation}
  \ln \tau_{kl,t}
  \;=\;
  \ln\bar\tau_{kl}
  + \lambda_{\phys}\ln D_{kl,t}
  + \eta_{kl} + \delta_t + \varepsilon_{kl,t},
\end{equation}
with IVs from weather shocks, policy timing, or upstream demand shifts
(Yang et al.\ 2020; Barwick et al.\ 2024).  Target: replace external
$\lambda_{\phys}=1.1$ with a Beijing-specific structural estimate and
s.e.

\subsection{Commuting gravity ( $\theta$ robustness)}

PPML or IV version of \eqref{eq:gravity_log} following Redding and Sturm
(2024), using straight-line distance or pre-period network features as
instruments.  Would complement, not replace, external $\theta=6.83$ unless
internal and external estimates converge.

\subsection{Sorting parameters $(\alpha,\beta)$}

Possible moments (not yet coded):
\begin{itemize}[leftmargin=2em]
  \item Cross-sectional variation in $(l^R,l^F)$ vs.\ model-implied
        market access under exogenous $\tau$ shocks;
  \item Second-difference orthogonality of $\Delta\ln\ubar_i$,
        $\Delta\ln\Abar_j$ to non-local transport improvements (Path
        Dependence \S3.3 style), adapted to tidal / lane shocks.
\end{itemize}
Current practice: $\alpha,\beta$ fixed; extensive sensitivity grid
(\S\ref{sec:sorting}).

\subsection{Land supply / amenity elasticities}

Not applicable in the current commuting-only v4 scaffold (no explicit
housing market block).  Deferred to extensions with floorspace data.

%====================================================================
\section{Results Registry and Reproducibility}\label{sec:registry}
%====================================================================

\begin{longtable}{@{}p{0.28\linewidth}p{0.34\linewidth}p{0.32\linewidth}@{}}
\toprule
\textbf{Analysis} & \textbf{Driver script} & \textbf{Output directory} \\
\midrule
\endhead
Stage2 baseline + CF &
  \texttt{run\_stage2\_gate5.py} &
  \texttt{stage2\_endogenous\_edgeproj\_v2\_lcc\_gate5\_bpr\_pos05\_v1/} \\
Tidal ST/LT diagnostics &
  \texttt{run\_tidal\_lane\_diagnostics\_gate5.py} &
  \texttt{tidal\_lane\_diagnostics\_...\_bpr\_pos05\_v1/} \\
Welfare / sorting decomposition &
  \texttt{run\_welfare\_sorting\_diagnostics\_gate5.py} &
  \texttt{welfare\_sorting\_diagnostics\_gate5\_bpr\_pos05\_v1/} \\
$\alpha,\beta$ sensitivity &
  \texttt{run\_alpha\_beta\_sensitivity\_bpr\_gate5.py} &
  \texttt{alpha\_beta\_sensitivity\_bpr\_pos05\_v1/} \\
Policy comparison P1--P7 &
  \texttt{run\_policy\_comparison\_gate5.py} &
  \texttt{policy\_comparison\_gate5\_bpr\_pos05\_v1/} \\
Exog vs.\ endog CF table &
  \texttt{run\_exog\_endog\_comparison\_gate5.py} &
  \texttt{exog\_endog\_comparison\_gate5\_bpr\_pos05\_v1/} \\
Fundamental maps &
  \texttt{run\_fundamental\_maps\_gate5.py} &
  \texttt{fundamental\_maps\_gate5\_bpr\_pos05\_v1/} \\
Full numerical summary &
  --- &
  \texttt{Documents/L2\_notes/model\_and\_results\_summary\_20260630.md} \\
\bottomrule
\end{longtable}

\paragraph{Canonical reproduction (mainline).}
\begin{verbatim}
cd data_work
python3 src/model/full_estimation/run_stage2_gate5.py \
  --pair-selection score_bpr_st --alpha 0.05 --beta -0.10 \
  --output-name stage2_endogenous_edgeproj_v2_lcc_gate5_bpr_pos05_v1
python3 src/model/full_estimation/run_tidal_lane_diagnostics_gate5.py \
  --only-top100 --pair-selection score_bpr_st
python3 src/model/full_estimation/run_welfare_sorting_diagnostics_gate5.py \
  --tracks pos05 --pair-selection score_bpr_st
\end{verbatim}

%====================================================================
\section{Parameter Summary Table}\label{sec:summary_table}
%====================================================================

\begin{center}
\begin{tabular}{@{}llll@{}}
\toprule
Parameter & Mainline value & Method & Override? \\
\midrule
$\theta$ & $6.83$ & AA (2022) Seattle / ARSW (2015) & External \\
$\alpha$ & $+0.05$ & Literature gross spillover (pos05) & Sensitivity \\
$\beta$ & $-0.10$ & AA (2022) Seattle & Sensitivity \\
$\lambda_{\phys}$ & $1.1$ & Barwick et al.\ (2024) mapping & $0.572$ band \\
$\lambda_{\AAcal}$ & $0.0714$ & $\delta_1/\theta$ (legacy code default) & Legacy form only \\
$(\ubar_i,\Abar_j)$ & inverted & Assumption~\ref{ass:obs_eq} & Fixed in CF \\
$\kappa$ & data-dependent & Spectral-radius bound & Algorithm \\
Pair selection & \texttt{score\_bpr\_st} & Mainline rule & Robustness \\
\bottomrule
\end{tabular}
\end{center}

%====================================================================
\section*{References}
\addcontentsline{toc}{section}{References}
%====================================================================

\begin{itemize}[leftmargin=2em]
  \item Ahlfeldt, G. M., Redding, S. J., Sturm, D. M., Wolf, N. (2015).
        ``The Economics of Density: Evidence from the Berlin Wall.''
        \emph{Econometrica}, 83(6), 2127--2189.
  \item Allen, T., Arkolakis, C. (2022). ``The Welfare Effects of
        Transportation Infrastructure Improvements.''
        \emph{Review of Economic Studies}, 89(6), 2911--2957.
  \item Barwick, P. J., Li, S., Lin, L. (2024). ``Traffic Congestion,
        Ambient Air Pollution, and Fuel Consumption: Evidence from
        Beijing.'' \emph{American Economic Review}.
  \item Redding, S. J., Sturm, D. M. (2024). ``Neighborhood Effects:
        Evidence from Wartime Destruction in London.'' NBER WP 32333.
  \item Yang, J., Purevjav, A. O., Li, S. (2020). ``The Marginal Cost
        of Traffic Congestion and Road Pricing: Evidence from a Natural
        Experiment in Beijing.'' \emph{AEJ: Economic Policy}, 12(1),
        418--453.
\end{itemize}

\end{document}
